\section{From Programs \& Types to Propositions \& Proofs}

The denotational soundness Theorem~\ref{thm:safety}
lets us interpret well typed programs as proofs of
propositions.

\NV{say that definition is a context that will be used later}
\mypara{``Definitions''}
A \emph{definition} $\defn$ is a sequence of reflected binders:
%
$$\defn \ ::= \ \bullet \spmid \erefb{x}{\gtyp}{e}{\defn}$$
%
A \emph{definition's environment} $\env(\defn)$ comprises
its binders and their reflected types:
%
\begin{align*}
\aenv(\bullet)                    \defeq & \emptyset \\
\aenv(\erefb{f}{\gtyp}{e}{\defn}) \defeq & (f, \exacttype{\gtyp}{e}),\ \env(\defn) \\
\end{align*}
%
A \emph{definition's substitution} $\sto(\defn)$ maps each binder
to its definition:
%
\begin{align*}
\sto(\bullet)                     \defeq & \emptysto \\
\sto(\erefb{f}{\gtyp}{e}{\defn})  \defeq & \extendsto{f}{\efix{f}\ e}{\sto(\defn)}
\end{align*}

\mypara{``Propositions''}
%
A \emph{proposition} is a type
%
$$\tbind{x_1}{\typ_1} \rightarrow \ldots
  \rightarrow \tbind{x_n}{\typ_n}
  \rightarrow \tref{v}{\tunit}{\ppn}$$
%
For brevity, we abbreviate propositions like the above to
%
$\tfun{\overline{x}}{\overline{\typ}}{\ttref{\ppn}}$
%
and we call $\ppn$ the \emph{proposition's refinement}.
%
For simplicity we assume that $\freevars{\typ_i} = \emptyset$.

\mypara{``Validity''}
%
\NV{add termination: proofs should provably terminates}

A proposition $\tfun{\overline{x}}{\overline{\typ}}{\ttref{\ppn}}$
is \emph{valid under} $\defn$ if
%
$$\forall \overline{w} \in \interp{\overline{\typ}}.\
  \evalsto{\applysub{\sto(\defn)}{\SUBST{\ppn}{\overline{x}}{\overline{w}}}}{\etrue}$$
%
That is, the proposition is valid if its refinement
evaluates to $\etrue$ for every (well typed)
interpretation for its parameters $\overline{x}$
under $\defn$.

\mypara{``Proofs''}
%
A binder $\bd$ \emph{proves} a proposition $\gtyp$ under $\defn$ if
$$\hastype{\emptyset}{\defn[\eletb{x}{\typ}{\bd}{\eunit}]}{\tunit}$$
%
That is, if the binder $\bd$ has the proposition's type $\gtyp$
under the definition $\defn$'s environment.

\begin{theorem}{[Proofs]} \label{thm:validity}
If $\bd$ proves $\typ$ under $\defn$ then $\typ$ is valid under $d$.
\end{theorem}

\begin{proof}
As $\bd$ proves $\typ$ under $\defn$, we have
%
\begin{align}
\hastype{\emptyset}{\defn[\eletb{x}{\typ}{\bd}{\eunit}]}{\tunit}
\label{pf:1} \\
%
\intertext{By Theorem~\ref{thm:safety} on \ref{pf:1} we get}
%
\sto(\defn) \in \interp{\env(\defn)} \label{pf:2}\\
%
\intertext{Furthermore,  by the typing rules \ref{pf:1}
implies $\hastype{\env(\defn)}{\bd}{\typ}$ and hence, via Theorem~\ref{thm:safety}}
%
\forall \sub \in \interp{\env(\defn)}.\ \applysub{\sub}{\bd} \in \interp{\applysub{\sub}{\typ}}
\label{pf:3} \\
\intertext{Together, \ref{pf:2} and \ref{pf:3} imply}
\applysub{\sto(\defn)}{\bd} \in \interp{\applysub{\sto(\defn)}{\typ}}
\label{pf:4}
\intertext{By the definition of type denotations, we have}
%
\interp{\applysub{\sto(\defn)}{\typ}}
  \defeq \{ f\ |\ \typ \mbox{ is valid under}\ \defn\}
  \label{pf:5}
\end{align}
By \ref{pf:4}, the above set is not empty, and hence $\typ$ is valid under $\defn$.
\end{proof}


%%%%%  \begin{definition}[Theorem]
%%%%%  Let $\aenv=\aenv(\prog)$ be the set of axioms for some program \prog.
%%%%%  Then
%%%%%  $\typ\defeq\tbind{x_1}{\typ_1} \rightarrow \dots \rightarrow \tbind{x_n}{\typ_n} \rightarrow \ttreft{v}{\btyp}{\refa}$
%%%%%  is a theorem if
%%%%%  $\freevars{\refa} \subseteq \{x_1, \dots, x_n\} \cup \domain{\aenv}$,
%%%%%  and for simplicity $\freevars{\typ_i} = \emptyset$.
%%%%%
%%%%%  Moreover, every expression $e'$ such that  \hastype{\aenv}{e'}{\typ}
%%%%%  provides a proof of the theorem $\typ$ with respect to the definitions in \prog.
%%%%%  \end{definition}
%%%%%  %
%%%%%  In other words, any expression $e'$ provides an evidence that the theorem expressed
%%%%%  by \typ holds.
%%%%%  The above definition is an instance of the Curry-Howard correspondence,
%%%%%  where programs as interpreted as proofs or the theorems expressed by their types.
%%%%%  %
%%%%%  %% Moreover, since the proof of the theorem is any expression $e'$ with type $\typ$
%%%%%  %% the definition states that our proof system is proof irrelevant.
%%%%%
%%%%%  \begin{theorem}
%%%%%  Let $\aenv=\aenv(\prog)$.
%%%%%  For every theorem
%%%%%  $\typ \equiv \tbind{x_1}{\typ_1} \rightarrow \dots \rightarrow \tbind{x_n}{\typ_n} \rightarrow \ttreft{v}{\btyp}{\refa}$
%%%%%  with a proof $e_\typ$ with respect to $p$,
%%%%%  then $\forall x_1\in\interp{\typ_1},\dots, x_n\in\interp{\typ_n}. \evalsto{\applysub{\Theta_\aenv(\prog)}{\refa}}{\etrue}$.
%%%%%  \end{theorem}
%%%%%  \begin{myproof}
%%%%%  Since \hastype{\aenv}{e_\typ}{\typ}
%%%%%  by Theorem~\ref{thm:safety} we have
%%%%%  $\forall \sub \in \interp{\aenv}.
%%%%%  \applysub{\sub}{e_\typ} \in \interp{\applysub{\sub}{\typ}}
%%%%%  $.
%%%%%  %
%%%%%  Since \prog type checks
%%%%%  we have $\Theta_\aenv(\prog)\in\interp{\aenv}$, thus
%%%%%  $\applysub{\Theta_\aenv(\prog)}{\refa} \in \interp{\applysub{\Theta_\aenv(\prog)}{\typ}}$.
%%%%%  %
%%%%%  By the Definition of Type Denotations we have
%%%%%  $\interp{\typ} =  \{f |\forall x_1\in\interp{\typ_1},\dots, x_n\in\interp{\typ_n}. \evalsto{\applysub{\Theta_\aenv(\prog)}{\refa}}{\etrue} \}$.
%%%%%  %
%%%%%  But $e_\typ$ is a witness that the set \interp{\typ} is not empty, thus
%%%%%  $\forall x_1\in\interp{\typ_1},\dots, x_n\in\interp{\typ_n}. \evalsto{\applysub{\Theta_\aenv(\prog)}{\refa}}{\etrue}$.
%%%%%  \end{myproof}

%% Theta_\aenv(\eletrec{f}{e}{\gtyp}{L}{\prog})
 %% & = (f, \efix{f}\ e),\aenv(p) \\
%% \Theta_\aenv(\elet{f}{e}{\gtyp}{L}{\prog})
 %% & = (f, e),\aenv(p) \\
%% \Theta_\aenv(\eletrecopt{f}{e}{\gtyp}{}{\prog})
 %% & = \Theta_\aenv(p) \\
%% \Theta_\aenv(e) &= \emptyset
%% \end{align*}

%% A proposition is \emph{well-formed} under $\defn$ if
%% $\freevars{\refa} \subseteq \{x_1, \dots, x_n\} \cup \domain{\aenv}$,


\mypara{Example: Fibonacci is increasing}
%
In \S~\ref{sec:overview} we verified that
under a definition $\defn$ that includes \fibname,
the term \fibincrname proves
$${\tfun{n}{\tnat}{\ttref{\fib{n} \leq \fib{(n+1)}}}}$$
%
Thus, by Theorem~\ref{thm:validity} we get
%that the \fib{}, as defined in $\defn$ is increasing
%
%$$
%\forall n\in\interp{\tnat}. \evalsto{\fib{n} \leq \fib{(n+1)}}{\etrue}
%$$
%Equivalently
$$
\forall n. \evalsto{0 \leq n}{\etrue} \Rightarrow \evalsto{\fib{n} \leq \fib{(n+1)}}{\etrue}
$$